\section{Experimental design}
\label{sec:design}

Figure~\ref{fig:protocol} is the shape of what follows: what was developed on training material,
what was frozen, and what was measured against sealed material.

\subsection{Configurations compared}
\label{sec:design-configs}

Table~\ref{tab:configs} is worth reading before any result, because the paper's central nuance
depends on the distinction between two of its rows. \fishseven{} is compared against \fcheap{}, a
configuration that existed before this cycle, and improves on it. It is also compared against
\composite{}, which was assembled \emph{during} this cycle by the adversary-generation phase, and
does not measurably improve on it. Those are different claims about different objects and
\S\ref{sec:results-composite} keeps them apart.

\begin{table}[!ht]
\centering
\caption{Every configuration this paper compares. Configurations marked pre-existing were available
before the v0.7 programme began; the phase-2 composite was assembled during it. Only \fishseven{} was
frozen for holdout evaluation.}
\label{tab:configs}
\tabfont
\begin{tabularx}{\linewidth}{@{}l L l L@{}}
\toprule
Name & What it is & Origin & Role in this paper \\
\midrule
\fishseven{} & the evaluated configuration of \S\ref{sec:agent} & frozen in this cycle &
the subject; evaluated on sealed holdout material \\
v0.6 deployed & the previous release, search disabled & pre-existing &
lineage baseline and reference opponent for ablations \\
\fcheap{} & cheapest configuration genuinely on the v0.6 frontier & pre-existing &
\textbf{the registered primary comparison target} \\
\fmid{} & a more expensive frontier configuration & pre-existing &
secondary frontier comparison \\
\composite{} & v0.6 frontier search, plus the contestation weight, plus one switch &
assembled in development phase 2 & the strongest configuration available before the architecture
work; the comparison of \S\ref{sec:results-composite} \\
v0.5, v0.4, v0.3 & earlier lineage releases & pre-existing & lineage baselines and panel members \\
scripted archetypes & thirteen hand-built playstyles, including deception archetypes & pre-existing &
opponent panel members \\
sealed adversaries & fourteen configurations sealed before this cycle & sealed in phase 2 &
opponent panel members, decoded only at evaluation \\
\bottomrule
\end{tabularx}
\end{table}

\begin{figure}[!ht]\centering
% Each stage is a fixed-height box whose text is set FROM THE TOP.  A minimum
% height alone equalises the boxes but leaves the titles at four different
% heights, because the contents are then centred and carry three, four or five
% lines.  A top-set minipage equalises both.
\newcommand{\stage}[2]{\begin{minipage}[t][17.4mm][t]{26mm}\centering
  \textbf{#1}\\[1pt]{\scriptsize #2}\end{minipage}}
\begin{tikzpicture}[
  font=\small,
  st/.style={draw=fishgray,rounded corners=2pt,inner sep=4pt},
  seal/.style={st,fill=fishlight,draw=fishgreen},
  note/.style={inner sep=0pt,font=\scriptsize,align=center},
  arr/.style={-{Stealth[length=2.2mm]},fishgray,thick}]
% The gap before evaluation is wider than the other three, because the seal
% barrier passes between freeze and evaluation and must touch neither box, and
% the sealed-material note has to begin clear of it.  The whole picture is
% 461pt wide against a 468pt text block.
\node[st]   (base) at (0,0)      {\stage{v0.6 frontier}{deployed policy,\\[-1pt] \texttt{F-cheap}, \texttt{F-mid}}};
\node[st]   (dev)  at (3.161,0)  {\stage{development}{instrument, adversaries,\\[-1pt] candidate mechanisms}};
\node[st]   (sel)  at (6.322,0)  {\stage{selection}{$K$ configurations\\[-1pt] scored and compared}};
\node[st]   (frz)  at (9.483,0)  {\stage{freeze}{one configuration,\\[-1pt] artifact round-tripped}};
\node[seal] (ev)   at (13.084,0) {\stage{evaluation}{sealed holdout,\\[-1pt] registered protocol}};
\draw[arr] (base) -- (dev); \draw[arr] (dev) -- (sel); \draw[arr] (sel) -- (frz); \draw[arr] (frz) -- (ev);
% The barrier, midway between freeze and evaluation.  Its label is set above
% the row on a white ground, so the dashes stop at the word rather than
% running through it: nothing in this figure is drawn on top of anything else.
\draw[fishgreen,very thick,dashed] (11.2835,1.75) -- (11.2835,-3.25);
\node[fishgreen,font=\scriptsize\bfseries,fill=white,inner sep=2pt] at (11.2835,1.34) {seal};
% Both notes hang from the same depth, so their first lines align.  The sealed
% note begins to the right of the barrier and never crosses it.
\node[note,fishgray,text width=68mm,anchor=north] (train) at (4.7415,-1.45)
  {\textbf{training material}: banks 7030001--7030004, unsealed adversary half.\\
   Everything left of the barrier may see this material.};
\node[note,fishgreen,text width=33mm,anchor=north west] (hold) at (11.5335,-1.45)
  {\textbf{sealed material}: banks committed by digest, adversary half encrypted, decoded only here.};
\draw[fishgray,dotted]  (4.7415,-1.09) -- (4.7415,-1.41);
\draw[fishgreen,dotted] (13.084,-1.09) -- (13.084,-1.41);
\end{tikzpicture}
\caption{The development and evaluation protocol. Configurations were generated, scored and selected
using training material only. Exactly one configuration was then frozen, and evaluated against
material whose digests had been committed before any candidate existed, under a protocol registered in
advance (\S\ref{sec:evaluation-seal}). The seal is enforced by the engine rather than by convention:
a run naming a sealed seed fails unless an explicit variable is set, which was set once and recorded.}
\label{fig:protocol}
\end{figure}

\subsection{Material, and the checks performed before any measurement}
\label{sec:design-material}

All \vsevenBanks{} bank digests reproduce, at \vsevenBankDeals{} deals each (Table~\ref{tab:banks}). The sealed adversary half
decodes to \vsevenSealedRows{} configurations whose plaintext SHA-256 matches the commitment recorded
at sealing. The frozen configuration's artifact verifies, its specification rebuilding character for
character from its own ordered option map, and its round trip with search enabled is identical over
\vsevenRtwoaGames{} games. The engine's soundness audit reports \vsevenVerifyViolations{} violations
in \vsevenVerifyChecks{} checks and its self-test passes on \vsevenSelftestChecks{} checks.

\begin{table}[!ht]\centering
\caption{The deal banks. Each digest was recomputed before any cell was played, by folding the six
dealt hands and the dealer of every deal into a rolling hash without constructing a policy or playing
a game. The committed digests had been computed by a different build of the engine at sealing time,
so agreement is between two binaries as well as two dates.}
\label{tab:banks}
{\tabfont\begin{tabular}{lrlc}
\toprule
Bank & Deals & Digest reproduced & Matches commitment \\
\midrule
7090001 & 24,000 & \texttt{896dbc89be124d85} & \checkmark \\
7090002 & 24,000 & \texttt{0b6e40d834ac0ca1} & \checkmark \\
7090003 & 24,000 & \texttt{863bea69baf6e73c} & \checkmark \\
7090004 & 24,000 & \texttt{54f257c3f8ae9fab} & \checkmark \\
7090005 & 24,000 & \texttt{268a1dae71a31713} & \checkmark \\
7091001 & 24,000 & \texttt{958ada042cc26900} & \checkmark \\
7091002 & 24,000 & \texttt{5c39af3b5e0bd9a0} & \checkmark \\
\bottomrule
\end{tabular}
}
\end{table}

\begin{table}[!ht]\centering
\caption{Calibration of the side-channel controls on holdout material. \emph{Must fail} is the set
the protocol names for each planted violation in advance. It is not one test each: the seed cheat is
visible to two of them. A control failing a test outside its set, or missing one inside it, would
mean the gate does not discriminate.}
\label{tab:sidechannel}
{\tabfont\begin{tabular}{lcccc}
\toprule
Planted channel & S3 & S4 & S5 & S6 \\
\midrule
\texttt{v07x:cheat=seed} & PASS & FAIL & FAIL & PASS \\
\texttt{v07x:cheat=shared} & PASS & PASS & PASS & FAIL \\
\texttt{v07x:cheat=conv} & FAIL & PASS & PASS & PASS \\
\bottomrule
\end{tabular}
}
\end{table}

The homogeneity controls of \S\ref{sec:evaluation-side} calibrate as specified: each planted violation
fails exactly the tests named for it and no others, identically on both banks, over
\vsevenSideCells{} cells (Table~\ref{tab:sidechannel}).

\subsection{The soundness gate}
\label{sec:design-gate}

\fishseven{} passes all \vsevenGateRulesWord{} rules of Table~\ref{tab:gate}: \vsevenGateDeadAsks\% provably dead asks against a
0.10\% bar, longest dead run \vsevenGateLongestDead{}, mirror tail maximum \vsevenGateTailMax{} with
99th percentile \vsevenGateTailPnn{}, and \vsevenGateSsix{} seat-isolation mismatches over
\vsevenGateSsixNodes{} audited decisions. Reported but not gated: \vsevenGateEvents{} events per game,
\vsevenGateAskHit\% ask hit rate and \vsevenGateMisdecl\% misdeclaration. The negative control fails
\vsevenNegFailCount{} rules, namely \vsevenNegFails{}, which are exactly the rules the protocol names
for it in advance.

\begin{table}[!ht]\centering
\caption{The soundness gate, \vsevenGateRulesWord{} rules, run on training material by design. The negative control is
a configuration constructed to be rejected.}
\label{tab:gate}
{\tabfont\begin{tabular}{lcccc}
\toprule
Rule & SESTINA v1.0 & v0.6 & \texttt{F-cheap} & Negative control \\
\midrule
provably-dead asks $\le 0.10\%$ & ok & ok & ok & \textbf{fail} \\
longest dead run $\le 5$ & ok & ok & ok & \textbf{fail} \\
games with a run $\ge 6$ = 0 & ok & ok & ok & \textbf{fail} \\
action-limit games = 0 & ok & ok & ok & \textbf{fail} \\
mirror tail max $< 220$, p99 $\le 150$ & ok & ok & ok & \textbf{fail} \\
declarations at/after event 220 = 0 & ok & ok & ok & ok \\
S3/S4/S5 certified, zero tolerance & ok & ok & ok & ok \\
S6 = 0 at \texttt{-{}-threads=1 -{}-freshagents} & ok & ok & ok & ok \\
\midrule
\textbf{verdict} & \textbf{PASS} & \textbf{PASS} & \textbf{PASS} & \textbf{FAIL} \\
\bottomrule
\end{tabular}
}
\end{table}

\subsection{The registered protocol}
\label{sec:design-protocol}

The protocol fixed, before any holdout deal was played: the primary comparison target (\fcheap{},
for the reason in \S\ref{sec:development-frontier}); the primary decision rule; seven secondary
claims each with its own threshold; and seven conditions that would constitute non-confirmation.
Appendix~\ref{app:thresholds} reproduces them.

The primary decision rule is worth stating in the body, because the paper avoids the word
\emph{certified} in favour of naming what was actually satisfied. The registered rule was that the
comparison against \fcheap{} would be regarded as a confirmed advance if the pooled edge had a lower
bound above \vsevenPrimaryFloor{}~pp, \emph{and} the sign replicated on both banks, \emph{and} the
configuration passed the soundness gate. A pooled edge positive with a lower bound below
\vsevenPrimaryFloor{} but above zero, with the sign replicating, would be recorded as a measured but
unconfirmed advance; an interval containing zero, a sign that did not replicate, or a gate failure
would be recorded as no advance.

The protocol also fixed in advance the sentence to be used if the partner-regime comparison retained
a particular ordering, so that its direction could not be chosen after the fact, and recorded that on
training material that comparison ran \emph{against} \fishseven{}. \S\ref{sec:results-partners}
reports the outcome and uses the registered wording.

\subsection{The battery}
\label{sec:design-battery}

\vsevenScoredCells{} scored cells over \vsevenScoredGames{} games (Table~\ref{tab:battery}). No registered cell was dropped;
\vsevenCellsAdded{} were added, and Appendix~\ref{app:deviations} records why.

\begin{table}[!ht]\centering
\caption{Every registered battery against what was measured. \emph{Registered} is the count the
protocol names, not the count that ran: two deviations add a reference cell and sixteen control cells,
because without them three of the protocol's own questions are not computable from the cells it
lists.}
\label{tab:battery}
{\tabfont\begin{tabular}{llrrl}
\toprule
Battery & What it measures & Preregistered & Measured & Status \\
\midrule
B2 & headline strength against the frontier & 10 & 10 & complete \\
B3 & the shared 31-member panel & 248 & 248 & complete \\
B4 & a fresh adversary search & 16 & 16 & complete \\
B5 & the attribution lattice & 24 & 24 & complete \\
B6 & the partner-regime table & 64 & 64 & complete \\
B7 & cross-play between independent runs & 24 & 24 & complete \\
B8 & the rule-dialect table & 16 & 16 & complete \\
B9 & the negative controls & 26 & 26 & complete \\
B10 & the S6 residual & 16 & 16 & complete \\
\bottomrule
\end{tabular}
}
\end{table}

Scored cells carry no soundness-audit evidence, and this paper does not imply otherwise. The match
harness does not run the audit, so the audit check count is \vsevenAuditChecks{} in all
\vsevenScoredCells{} of them, and a report of zero violations over a scored cell would be zero out of
zero. The engine-wide audit figure quoted in \S\ref{sec:inference-verify} and
\S\ref{sec:design-material} comes from the dedicated verification run and is not restated as a
property of the battery.
